\subsection*{Chapter 1: Euclidean Geometry and Cartesian Coordinates}

The first chapter is designed to remove a common but misleading impression: that
geometry begins when axes and ordered pairs are drawn. The Euclidean plane already
contains points, lines, circles, angles, congruence, perpendicularity,
parallelism, and distance. Coordinates are an additional structure placed on
that plane. They are a method of naming and calculating with geometric objects,
not the source of the objects themselves.

The chapter therefore moves from geometry to arithmetic rather than beginning
with arithmetic and pretending that the geometry has already been explained. Its
central intention is to teach the student to distinguish an object from a
representation of that object. A point is not an ordered pair by nature; an angle
is not a number by nature; a length is not created by writing a decimal. Numbers
become attached to these objects only after choices of origin, direction, unit,
and measurement convention have been made.

\subsubsection*{The Euclidean geometry available before coordinates}

This section establishes the world in which the later coordinate construction
will take place. The distinction between Euclid's postulates,
straightedge-and-compass operations, and derived constructions is essential. A
postulate describes the geometric setting being assumed. An operational rule
states what elementary action may be performed. A derived construction must be
built from those actions and justified. The student should therefore learn that
``we can draw it'' and ``we have proved that the construction has the required
property'' are different claims.

The perpendicular bisector, midpoint, perpendicular through a point, angle
bisector, and parallel through a point are deliberately constructed rather than
silently treated as primitive tools. They show how complex geometric abilities
grow from a small set of permitted actions. The student should become able to
read a construction as an argument: every new point comes from an allowed
intersection, every equality of lengths comes from a circle or congruence, and
every final claim must follow from facts already established.

Parallel lines then lead to proportional segments and similarity. This prepares
a second important idea: a quantity may be independent of the size of the figure
used to represent it. Sine and cosine arise as ratios that remain unchanged among
similar right triangles. The identity
\[
\cos^2\alpha+\sin^2\alpha=1
\]
is not presented as a detached trigonometric rule, but as the Pythagorean theorem
written in ratio form.

The discussion of angle measure makes a further distinction. An angle exists as
a geometric region before degrees or radians are assigned to it. Radian measure
arises from the requirement that the number attached to an angle should not
depend on the radius of the circle used to measure it. The ratio of arc length to
radius satisfies that demand. The law of cosines then shows that an angle may be
recovered from three Euclidean lengths before vectors or Cartesian coordinates
have been introduced.

By the end of this section, the student should be able to distinguish assumptions,
allowed operations, constructions, and theorems; justify basic Euclidean
constructions; explain why trigonometric ratios are well defined; distinguish an
angle from its measure; and recognise that numerical formulas can arise from
purely geometric demands.

\subsubsection*{Constructing Cartesian coordinates from Euclidean geometry}

The second section asks what additional choices are required to turn the already
existing Euclidean plane into a Cartesian coordinate plane. The construction is
intentionally gradual. We choose a first line, select an origin, construct a
perpendicular axis, choose positive and negative directions, transfer a common
unit length, refine the scale, and finally assign coordinates by orthogonal
projection.

Each step removes an ambiguity. A line alone gives no origin. An origin gives no
orientation. Two perpendicular axes give no scale. Integer marks do not yet give
a complete real-number coordinate system. The continuity assumption is therefore
stated openly: finite straightedge-and-compass constructions produce many
specific points, but identifying every point of an axis with exactly one real
number requires an additional completeness principle. The student should learn
that mathematical structures often depend on choices and assumptions that are
hidden when only the finished notation is shown.

Orthogonal projection then turns a point into two signed numerical labels. The
coordinate pair records where the point lies relative to the chosen axes, not an
intrinsic identity of the point. Finally, the Euclidean distance formula is
derived from the Pythagorean theorem. This closes the chapter by showing how a
coordinate calculation inherits its meaning from the geometry that existed
before coordinates.

The intended competence is not merely the ability to plot a pair $(x,y)$. The
student should be able to reconstruct conceptually what a Cartesian system
contains, explain the role of each choice, use projections to interpret
coordinates, and derive rather than merely quote the distance formula. Most
importantly, the student should understand that coordinates translate geometry
into arithmetic while leaving the underlying geometric object unchanged.
